\section{前沿：四组无法被彻底消除的张力}
\label{sec:frontiers}

编译技术的前沿并不是一列等待安装的新功能。验证编译、超优化、PGO、JIT、异构调度和多层 IR 之所以持续演进，是因为它们分别在几组不可同时取满的目标之间移动边界。设候选编译方案为 $x$，我们可以用一个帕累托（Pareto）问题概括这些选择：
\begin{equation}
  \underset{x\in\mathcal{X},\;x\models\mathcal{S}}{\operatorname{Pareto\!\!-min}}
  \bigl(
    C_{\mathrm{run}}(x),
    C_{\mathrm{compile}}(x),
    C_{\mathrm{memory}}(x),
    C_{\mathrm{maint}}(x),
    R_{\mathrm{semantic}}(x)
  \bigr).
  \label{eq:frontier-pareto}
\end{equation}
$\mathcal{X}$ 是可考虑的实现与调度，$\mathcal{S}$ 是必须满足的语义和平台约束；各分量分别代表运行、编译、内存、维护成本与残余语义风险。现实系统很少存在支配所有其他方案的唯一最优点。所谓进步，常常是扩展可行集、降低某项成本，或让原本隐含的风险可以度量。

\subsection{保证范围与适用覆盖}

验证编译器希望对一族输入程序证明语义保持。若源程序为 $S$、目标程序为 $T$，其基本形态是行为细化（behavioral refinement）
\begin{equation}
  T=C(S)\Longrightarrow \operatorname{Beh}(T)\subseteq\operatorname{Beh}(S),
  \label{eq:frontier-refinement}
\end{equation}
但 $\operatorname{Beh}$ 如何处理未定义行为、I/O、发散和并发，本身就是定理的一部分。CompCert 通过较受控的语言与流水线获得跨多道 Pass 的机器证明，并缩小可信计算基（trusted computing base）\cite{adv-compcert-cacm,adv-compcert-site,adv-compcert-tcb}；Alive2 则选择翻译验证（translation validation），逐次检查 LLVM IR 优化前后的细化关系，从而更容易接入快速演化的工业编译器，但受到操作覆盖、循环界和求解预算限制\cite{adv-alive2}。

二者揭示的不是“证明优于测试”这样简单的层级，而是保证范围与生态覆盖之间的交换。模型越完整、可信计算基越小，支持新语言特征和工具链组件的成本通常越高；覆盖越广、演化越快，就越需要局部验证器、测试和工程约束共同承担信任。

本文案例中的对应决策是把关键许可压缩为一个容易检查的不变量：$0\le k\le8$ 且每项限幅到 $[-8,12]$，因此整数归约不会溢出。我们不需要证明整条工具链，便能说明为什么加法重结合在这个域内合法；若未来去掉长度上界，则必须撤回这项优化许可。前沿方向是让这种局部事实更容易穿过 IR，被重写系统和验证器共同消费，而不是把每个课程例子都包装成端到端形式化工程。

\subsection{搜索空间与编译预算}

经典优化器常以贪心次序应用规则，速度可控，却可能因为早期选择而错过更好的组合。多面体模型（polyhedral model）把仿射循环的迭代域、访问和依赖变成整数约束，使交换、融合、分块与并行化成为调度搜索\cite{adv-pluto,adv-polly}；超优化（superoptimization）直接搜索短程序空间\cite{adv-stoke,adv-souper}；等式饱和（equality saturation）用等价图（e-graph）同时保存许多等价表达式，最后按成本抽取，从而推迟破坏性的规则次序\cite{adv-egg,adv-cranelift-egraph}。

候选越多并不自动意味着结果越好。搜索需要燃料上限，等价规则需要适用前提，成本模型还可能把指令数较少误判为真实机器上更快。工业系统因而倾向于缩小搜索域：只对无副作用区域饱和，只综合短整数片段，或离线发现规则后交给独立检查器。它们没有消灭组合爆炸，只是把预算花在收益密度最高的地方。

截断点积把这项张力变得具体。逐项的 \code{min/max}、比较与选择可以在受限 e-图中选择分支式或无分支式实现；外层归约可比较标量链、向量水平归约与展开树。但候选集合必须携带前节的不溢出条件，而且优化目标要针对 RV64GC 或实际 NPU，而不是抽象“操作数最少”。我们的设计决策因此是：在结构化 IR 中延迟循环形态，保留若干调度选择；进入后端后收窄空间，以目标成本和编译时限作出确定选择。

\subsection{静态知识与运行时知识}

静态编译器知道类型、控制依赖和目标 ISA，却不知道部署时最常出现的输入、真正的热路径或动态类型。配置文件引导优化（profile-guided optimization, PGO）把运行剖面带回下一次构建，BOLT 在最终二进制布局已经可见时重排代码，机器学习引导优化（machine-learning-guided optimization）尝试从历史程序学习复杂启发式\cite{adv-clang-pgo,adv-autofdo,adv-bolt,adv-mlgo-paper,adv-mlgo-doc}。即时编译（just-in-time compilation, JIT）则把决策继续推迟到运行期，以编译开销、预热和代码缓存换取真实类型与热度；分层运行时在快速基线代码与优化代码之间进一步折中\cite{adv-art-jit,adv-orc}。

这里没有脱离工作负载的“最佳编译时刻”。若截断点积只在短命进程中调用一次，JIT 的准备成本不可能由八次乘法偿还；若它成为长寿命服务中的热点，且 $k$ 几乎总为 $8$，运行时版本化可以把边界处理移出热路径。可是一旦输入分布改变，旧 profile 和学习模型就会成为过期的编译输入。静态证明决定候选是否允许，运行信息决定候选是否值得；前者不能由热点数据替代，后者也不能从语义定理推导。

本案例选择预先编译（ahead-of-time compilation, AOT）路径，是因为目标是解释从源程序到进程映像的组成，而非追求一个八元素内核的峰值性能。这个选择本身体现了式~\eqref{eq:frontier-pareto}：牺牲运行时特化机会，换取更小的系统状态、更直接的 ABI 和更清晰的观察路径。若场景改变，合理答案也应改变，而不是把 AOT 或 JIT 升格为原则。

\subsection{抽象能力与互操作边界}

高层抽象保存形状、归约和布局，使融合、自动调度与跨设备映射成为可能；低层接口稳定、工具丰富，也更容易同既有运行时、链接器和调试器互操作。MLIR、TVM、Halide 与 Tiramisu 等系统从不同角度分离计算含义与调度选择\cite{adv-mlir-cgo,adv-tvm,adv-halide,adv-tiramisu}。代价是每增加一种方言、布局或调度层，就增加一组语义接口、转换规则与失败诊断。抽象若没有明确的消去边界，最终只会把复杂性推给下游。

第~\ref{sec:mlir} 节中的点积路线因此做了两次有意的收缩。第一次在结构化层保留映射--限幅--归约，使 CPU 向量化或 NPU 分块仍有选择；第二次在选定目标后落实 \code{memref}、地址空间、ABI 和设备入口，使剩余系统能够依赖普通链接与运行时协议。Ascend 的官方 \code{VecAdd} 位于第二次收缩之后：它清楚展示 GM--UB 搬运和 HIVM 向量操作，却不承担高层归约调度。把这条边界说清，比画一条貌似全自动的端到端箭头更有价值。

这种张力也解释了为何全程序优化通常采用摘要与选择性导入，而不是永久维持一个巨大全局 IR。ThinLTO 用模块摘要换取跨翻译单元可见性，同时保留并行后端与缓存能力\cite{adv-thinlto-paper,adv-thinlto-doc}；Transform dialect 把部分调度策略显式表示出来，却仍需目标成本和可诊断的失败语义\cite{adv-transform-dialect}。成功的抽象不是永不泄漏，而是在泄漏发生时拥有窄而明确的接口。

\subsection{共同方向：使选择及其条件显式化}

四组张力最终汇合到同一个设计观：编译器的职责不是找到脱离条件的“最佳程序”，而是在可承担的知识、时间和信任预算下，选择一个满足语义约束的实现。形式化方法让允许什么更清楚；搜索方法扩大可以选择什么；profile 与 JIT 改变何时知道什么；多层 IR 决定哪些信息能活到选择发生的时刻。

对贯穿全文的截断点积，我们可以据此读出四个相互连接的决定：用范围不变量许可归约变换；用结构化表示保留有限的调度空间；因程序短小而选择 AOT，不引入无法摊销的动态状态；在 LLVM、ELF 和设备 ABI 边界主动消去高层抽象。换一个数值域、工作负载或硬件，这些决定都可能变化，但推理框架不变。

这也是“深入理解编译系统”比罗列 Pass 更重要的地方。编译器的每一层都在交换信息：丢掉一种自由，换来一种确定性；暴露一种机器事实，失去一种高层对称；扩大一种搜索能力，支付一种验证与维护成本。前沿研究真正改变设计空间之处，不是又增加一个术语，而是让这些交换可表达、可组合，并让工程师知道自己为何停在帕累托前沿上的这一点。
